\begin{statementbox}

	\textbf{\textcolor{CStatement}{条件}}
	\begin{description}[style=nextline, leftmargin=2.8em, labelsep=0.5em, itemsep=0.35em]
		\item[\textbf{\textcolor{CStatement}{条件 1（正弦、余弦公式）：}}]
		      正弦二倍角公式、余弦二倍角公式对任意实数 $x$ 都成立。

		\item[\textbf{\textcolor{CStatement}{条件 2（正切公式定义域）：}}]
		      若使用
		      \[
			      \tan 2x=\frac{2\tan x}{1-\tan^2 x},
		      \]
		      则需使右端有意义，即
		      \[
			      \cos x\ne 0,\qquad 1-\tan^2 x\ne 0.
		      \]
		      等价于
		      \[
			      x\ne \frac{\pi}{2}+k\pi,\qquad
			      x\ne \frac{\pi}{4}+\frac{k\pi}{2}\quad(k\in\mathbb Z).
		      \]
		      特别注意：$\tan 2x$ 本身存在只要求 $\cos 2x\ne 0$；这里额外要求 $\cos x\ne0$，是为了保证右端的 $\tan x$ 有意义。
	\end{description}

	\textbf{\textcolor{CStatement}{结论}}
	\begin{description}[style=nextline, leftmargin=2.8em, labelsep=0.5em, itemsep=0.45em]
		\item[\textbf{\textcolor{CStatement}{结论一（正弦二倍角）：}}]
		      \textbf{\textcolor{CStatement}{直观描述：}} 正弦函数的二倍角等于两倍正弦与余弦的乘积。
		      \[
			      \sin 2x = 2\sin x \cos x.
		      \]

		\item[\textbf{\textcolor{CStatement}{结论二（余弦二倍角）：}}]
		      \textbf{\textcolor{CStatement}{直观描述：}} 余弦函数的二倍角可表示为余弦平方减正弦平方等多种形式。
		      \[
			      \cos 2x = \cos^2 x - \sin^2 x,
		      \]
		      \[
			      \cos 2x = 1 - 2\sin^2 x,
		      \]
		      \[
			      \cos 2x = 2\cos^2 x - 1.
		      \]

		\item[\textbf{\textcolor{CStatement}{结论三（正切二倍角）：}}]
		      \textbf{\textcolor{CStatement}{直观描述：}} 正切函数的二倍角可用 $\tan x$ 表示为分式形式，但必须注意定义域。
		      \[
			      \tan 2x = \frac{2\tan x}{1 - \tan^2 x}.
		      \]
		      其中要求
		      \[
			      \cos x\ne0,\qquad 1-\tan^2x\ne0.
		      \]

		\item[\textbf{\textcolor{CStatement}{常见变形（降幂公式）：}}]
		      \textbf{\textcolor{CStatement}{直观描述：}} 利用余弦二倍角变形可得到正弦、余弦的降幂公式。
		      \[
			      \sin^2 x = \frac{1 - \cos 2x}{2},
		      \]
		      \[
			      \cos^2 x = \frac{1 + \cos 2x}{2}.
		      \]
	\end{description}

\end{statementbox}